% chapters/introduction.tex

\chapter{Introduction}{\label{ch:intro}}

\section{Background}

In modern academic institutions, efficient management of administrative and academic processes is
critical for ensuring smooth operation and high quality learning outcomes. However, many
institutions still rely on partially manual or loosely integrated systems for managing core activities
such as attendance tracking, scheduling, feedback collection, and academic reporting.

Traditional attendance systems often involve manual roll calls or basic digital entry, both of which are
prone to errors, proxy attendance, and inefficiencies. Similarly, scheduling is frequently handled
through static timetables with limited flexibility, leading to conflicts and lack of real time updates.
Feedback collection mechanisms, where implemented, are often inconsistent and fail to provide
actionable insights due to low participation or poor structure.

Furthermore, the absence of centralized systems results in fragmented data across different
processes. This makes it difficult for administrators and faculty to obtain a comprehensive view of
student engagement, faculty workload, and overall academic performance. As institutions scale,
these inefficiencies become more pronounced, highlighting the need for an integrated and
automated solution.

The CiPD (Centre for Intelligent Product Development) program at IIIT Delhi operates in a dynamic
and fast paced academic environment, where multiple sessions, instructors, and student interactions
occur daily. This setting demands a system that not only manages operations efficiently but also
provides real time insights and automation to support decision making.

\section{Problem Statement}

The primary problem addressed in this project is the lack of a unified, scalable, and automated
system to manage academic operations within the CiPD program.

Specifically, the following challenges were identified:

\begin{itemize}
    \item \textbf{Inefficient Attendance Tracking:}
    Existing attendance methods are either manual or semi digital, making them susceptible to
    inaccuracies, proxy attendance, and delayed processing.

    \item \textbf{Lack of Real Time Scheduling Management:}
    Scheduling of sessions is not dynamically managed, leading to potential conflicts, lack of
    visibility, and difficulty in handling changes.

    \item \textbf{Inadequate Feedback Collection and Analysis:}
    Feedback systems, if present, are not tightly integrated with sessions, resulting in low
    response rates and limited actionable insights.

    \item \textbf{Fragmented Data and Limited Analytics:}
    Data related to attendance, feedback, and academic activities is not consolidated, restricting
    the ability to perform meaningful analysis and generate reports.

    \item \textbf{Absence of Automation in Academic Workflows:}
    Key workflows such as attendance processing, feedback rollout, and notifications are not
    automated, increasing administrative overhead.
\end{itemize}

These challenges collectively reduce operational efficiency, limit data driven decision making, and
affect the overall academic experience for students, faculty, and administrators.

The objective is to design a scalable prototype that demonstrates the feasibility and effectiveness of
an integrated academic ERP system, with potential for further expansion and enhancement in future
iterations.
